\documentclass[12pt]{book}
\usepackage[utf8]{inputenc}
\usepackage[T1]{fontenc}
\usepackage{amsmath,amssymb,amsfonts}
\usepackage{mathrsfs}

% Calligraphic for categories and sheaves
\newcommand{\cat}[1]{\mathcal{#1}}
\newcommand{\sheaf}[1]{\mathcal{#1}}

% Standard operators
\DeclareMathOperator{\Hom}{Hom}
\DeclareMathOperator{\Ext}{Ext}
\DeclareMathOperator{\id}{id}
\DeclareMathOperator{\dimk}{dim_k}
\DeclareMathOperator{\Rfst}{\mathbf{R}f_*}
\DeclareMathOperator{\Lfst}{\mathbf{L}f_*}
\DeclareMathOperator{\RHom}{\mathbf{R}\mathcal{H}\textit{om}}

\begin{document}

\setcounter{page}{10}

(iv) Recently P. Deligne has shown (unpublished) that for the category of noetherian preschemes, one can construct \(f^{!}\) and \(\operatorname{Tr}_{f}\) satisfying (a1), (b1), and (c), working with \(F\in D(\cat{QCoh}(X))\) and \(G\in D^{+}(\cat{QCoh}(Y))\), the derived categories of the categories of quasi-coherent sheaves on \(X\) and \(Y\), respectively.

The principal difficulty has been the lack of a suitable local, then glue, construction of the functor \(f^{!}\). Our procedure is to define for a finite or a smooth morphism, we have it given to us, by (a2) and (a3) [III \S]. Thus by composition we can obtain it for any morphism which can be factored into a finite morphism followed by a smooth morphism [III \S]. However, the derived category is not a local object, so we cannot glue these local determinations of \(f^{!}\) to obtain a global one. We resort to a clumsy, round-about procedure of defining \(f^{!}\) for a special class of complexes, called residual complexes [VI \S], and then pulling ourselves up by our bootstraps to get it for arbitrary complexes [VII \S]. But our result has the unpleasant hypotheses of (ii) above.

Deligne's construction of \(f^{!}\) is entirely different, and proceeds essentially by representing the functor (for \(F\in D(X)\))
\[
F \longmapsto \RHom_{Y}\bigl(\Rfst F, G\bigr).
\]
This approach has the advantage of giving the duality theorem immediately. He also has a method for calculating \(f^{!}G\) locally on \(X\). However, it is not immediately clear from his construction that the properties (a2), (a3) and (b2), (b3) hold. In fact, their proof will probably require some knowledge of the duality theorem as we have proved it here. At least one can hope that some combination of the two approaches will provide substantial simplifications of the theory at a later date.

\setcounter{page}{12}

A second difficulty, which lurks on the fringes of these notes, is that the derived category seems to be a little too big when it comes to unbounded complexes.
For example, one can have two unequal morphisms \(f,g\colon X\to Y\) in \(D^{+}(\cat{A})\) (where \(\cat{A}\) is an abelian category) such that their restrictions to each truncation of the complex \(X\) to a bounded complex are equal. This gives rise to some trouble with unbounded complexes (see the problem after [II 5.7], the boundedness hypotheses in [IV 3.1] and [VII 4.3a], and the remark following [VI 1.1c]). Perhaps one will have to replace \(D^{+}\) by the categories \(\operatorname{ind}\,D^{b}\) and \(\operatorname{pro}\,D^{b}\), which however may not be triangulated categories.

\setcounter{page}{13}

Thirdly, some discussion of our noetherian hypotheses is in order. In the present state of the theory, the noetherian hypotheses are well entrenched. We have used them in [II \S] for the structure of the injective objects in the category of sheaves on a locally noetherian prescheme, and its consequence that \(D^{+}(\cat{QCoh}(X))\to D_{\mathrm{qc}}^{+}(X)\) is an equivalence of categories. We have used them in the construction of the trace map for projective space [III \S], in the construction of \(f^{!}\) for a finite morphism [III \S], in the whole theory of dualizing complexes [V \S2], in the definition of residual complexes [VI \S], and so forth. Our finite Krull dimension hypotheses are often needed only to make possible the definition of \(\Rfst\) for unbounded complexes --- a problem which will disappear when the second difficulty above is solved, and the relation between \(D(\cat{QCoh}(X))\) and \(D_{\mathrm{qc}}(X)\) is better understood. It is certainly clear that our methods of proof rely heavily on noetherian hypotheses. I expect, however, that once a suitable statement is obtained, e.g., in case (iii) above, one could expect to prove the theorem without noetherian hypotheses, by reducing to the noetherian case. More satisfactory, of course, would be a treatment where noetherian hypotheses on the base were eliminated from the proofs as well as from the statements. When that is achieved, it will be reasonable to state the duality theorem for a proper morphism of ringed topos, whose fibres are noetherian schemes\dots At the present, however, this must remain a dream of things to come.

Now we will give the reader a brief description of the organization of these notes.

\setcounter{page}{14}

Chapter I gives the language of derived categories, which is used continually in the sequel. Sections 1--5, containing the definition of derived categories and derived functors, are essential, while sections 6 and 7 are refinements which are needed in proofs later on. This chapter is a self-contained treatment of the subject, and makes no reference to algebraic geometry, so can be used independently as an introduction to the notion of derived category. Sections 1--6 (except for the notion of localizing subcategory and the corresponding categories \(K_{A'}(\cat{A})\), \(D_{A'}(\cat{A})\), etc.) are taken almost without change from notes of Verdier, and should appear in his thesis [18].

Chapter II is a fairly systematic treatment of the applications of the language of derived categories to the category of sheaves on a prescheme. We consider the functors \(\Gamma\), \(f_{*}\), \(\mathcal{H}\textit{om}\), \(\otimes\), \(f^{*}\), their derived functors, and relations between these derived functors, such as associativity formulae. There is really no new mathematics involved here, since it is merely a translation into a new language of known results. However, some care is taken to see what hypotheses are needed. Only section 7 is notable new material. Here we give the structure of injective \(\sheaf{O}_X\)-modules on a locally noetherian prescheme \(X\), showing that they are all direct sums of certain indecomposable injectives \(J(x,x')\), for pairs of points \(x\) specializing to \(x'\) of \(X\). This extends results of Matlis [13] and Gabriel [5] for the case of quasi-coherent sheaves.

\setcounter{page}{15}

Chapter III contains everything we can say about duality for projective morphisms, and if the reader cares only for them, he may stop at the end of this chapter. However, it should be noted that the following chapters, besides giving us the tools for the proof of the general duality theorem in Chapter VII, also give much new insight into the nature of duality for a projective morphism.

There are so many situations in which we have a functor \(f^{!}\) like the one mentioned in the ideal theorem above, that we use different notations for them: \(f^{*}\) for a smooth morphism in section 2, \(f^{p}\) for a finite morphism in section 6, and \(f^{!}\) for an embeddable morphism in section 8. Later we will also have \(f^{Y}\), \(f^{Z}\) [VI \S2] and \(f^{\Delta}\) [VI \S3] for residual complexes.

In sections 3, 4, and 5 we recall the explicit calculations of cohomology for projective space, and give the old duality for projective space in the new language of derived categories. Sections 6 and 7 give the corresponding formalism for finite morphisms, and its relation to the case of smooth morphisms, so that in sections 8, 10, and 11 we can prove the ideal duality theorem for projective morphisms. Section 9 gives the formalism of a residue symbol which generalizes the classical residue of a differential on a curve. Even over the complex numbers, this important concept of residue for varieties of dimension greater than one was not known before.

\setcounter{page}{16}

In Chapter IV we study ``local cohomology'', or cohomology with restricted supports, of abelian sheaves on a locally noetherian topological space, generalizing results of [LC \S]. In particular, we discuss various cohomological properties which a sheaf or complex of sheaves may have with respect to certain families of supports. This gives rise to the notions of depth, Cousin complex, Cohen--Macaulay complex, and Gorenstein complex. The results of this chapter are independent of all other chapters of these notes, and so may be of use elsewhere, although their only application for the moment is to the theory of duality on preschemes.

Chapter V discusses dualizing complexes (read section 0 to find out what one is) and gives a duality theorem for modules over a local ring, generalizing results of [LC \S]. In particular, the dualizing functor
\[
\underline{D} = \RHom(\,\cdot\,, R^\cdot)
\]
treated here will be useful for our bootstrap operation (construction \(f^{!}\)) later, because it interchanges \(\otimes\) and \(\Hom\), \(f^{!}\) and \(\Lfst\), and commutes with \(\Rfst\) (duality theorem).

In Chapter VI we prepare for the final duality theorem by giving the construction of \(f^{!}\) and \(\operatorname{Tr}_{f}\) for residual complexes.
This is accomplished by a delicate glueing procedure which is the most difficult part of the theory, so we have given it in some detail. Perhaps some day this type of construction will be done more elegantly using the language of fibred categories and results of Giraud's thesis [6].

\setcounter{page}{17}

Chapter VII contains two main results. The first is the residue theorem, which generalizes the classical theorem that the sum of the residues of a differential on a curve is zero, and which is proved by reduction to that case. The second is the proof of the duality theorem for a proper morphism, which, now that all the functorial machinery has been set up, is little more than putting together the pieces of a jigsaw puzzle.

It remains to give the reader some perspective by listing some topics which have a logical place in these notes, but which are not here.

1. The cohomology class associated to a cycle. Let \(X\) be proper and smooth of dimension \(n\) over a field \(k\). In [8, \S4] it is shown how to associate to each non-singular subvariety \(Y\) of codimension \(p\) of \(X\), a cohomology class
\[
P_{X}(Y) \in H^{p}\bigl(X,\Omega_{X/k}^{p}\bigr).
\]

\setcounter{page}{18}

Now this can be done for an arbitrary subvariety of \(X\), using the remarks in [9]. One defines \(\omega_{Y}=H^{-n+p}(g^! k)\), where \(g\colon Y\to \operatorname{Spec}k\) is the projection. Then there is a canonical element
\[
\eta \in \Ext_{\sheaf{O}_{Y}}^{-n+p}\bigl(\omega_{Y}, g^{!}k\bigr).
\]
One defines a natural map
\[
\Omega_{Y/k}^{n-p} \to \omega_{Y},
\]
which, together with \(\eta\) and the construction of [8, \S4], gives the cohomology class \(P_{X}(Y)\). One proves the fundamental theorem that formation of the cohomology class associated to a cycle takes an intersection of cycles into the cup-product of their cohomology classes.

2. The theory of Poincaré duality and the Gysin homomorphism can be developed as in [8, \S].

3. A Lefschetz--Verdier fixed point formula for coherent sheaves on a scheme proper over a field, to generalize [21, Thm.~2]. In particular, the determination of the local contribution at a non-simple fixed point presents an interesting topic for future investigation.

4. The still lacking theory of duality for complexes with differential operators as boundary operator, its ties with ordinary singular homology theory and with vector bundles with integrable connections, as suggested in [22].

\medskip
A remark on references: Theorem 5.1 of Chapter III is referred to as ``Theorem 5.1'' in Chapter III, and as [III 5.1] elsewhere. References to the bibliography at the end are given by square brackets with an arabic numeral or some capital letters e.g., [14,(31.I)] or [EGA III 2.1.12].

\end{document}